\textbf{Abstract:} \hspace{0.3cm} The OWASP Mobile Application Security project is the industry standard when it comes to android mobile application security. It comprises all aspects of mobile application security, from documenting weaknesses to providing verification standards. Part of it is the Mobile Application Security Testing Guide (MASTG), which gives detailed instructions on how to thoroughly test mobile applications. Complementing the MASTG are OWASP MASTG reference apps. These apps serve as a source of learning by including intentionally built-in weaknesses that can be found and exploited using the MASTG, teaching mobile developers what to avoid in their own applications. They can also be used as a dataset for evaluating and benchmarking the effectiveness of different static and dynamic application security testing tools.

However, many of the existing reference apps are limited in scope and have become outdated due to a lack of maintenance. This limits their function as both a learning resource for mobile developers and a dataset for evaluating security testing tools. This BA thesis aims to evaluate the state of the art of the MASTG reference apps and conduct a literature review on the relevance and usage of the MASTG and related OWASP fields in the academic world. Based on the findings, new MASTG reference apps are developed. These reference apps include both dynamic and static vulnerabilities and document each vulnerability in detail. The documentation entails a mapping of the implemented vulnerabilities to their respective OWASP counterparts and describes where the vulnerabilities can be found in the code and how they can be exploited. 

The resulting open-source OWASP MASTG reference applications can be used both as a resource for teaching mobile application security and for benchmarking security testing tools.